\scanpage{307}
ont même fibre générique, il s'ensuit que cette dimension est la même pour $G$ et $H$, donc $u_s:G_s\quad H_{\deleted{\unsure{g}}\inserted{s}}$ est un \uline{monomorphisme} de groupes algébriques \uline{de même dimension}. On en conclut facilement que $G_s$ est ouvert dans $H_s$, et en fait est ensemblistement une \deleted{\unsure{réun}} réunion de composantes connexes de $H_s$ (on est ramené au cas où le corps de base est algébriquement clos, et $G$ et $H$ réduits, donc simples sur $k$, où c'est immédiat ...). Donc $H_s-G_s$ est fermé dans $H_s$, donc dans $H$, donc son complémentaire $H'$ dans $H$ est ouvert, et c'est évidemment un ouvert stable \unsure{pourla|pour la} loi de groupe de $H$. Donc, \deleted{pour prouver que} \deleted{\unsure{$u$ est}} quitte à remplacer $H$ par $H''$, on peut, pour prouver que $u$ est une immersion, se ramener au cas où, en plus des hypothèses précédentes, on suppose $u$ \uline{bijectif}, \deleted{\unsure{en part}} i.e. $u_s:G_s\quad H_s$ \uline{bijectif}. On est donc en tous cas ramené à prouver que $u$ \inserted{ou encore $u_s$} est un isomorphisme, moyennant le cas échéant l'hypothèse de bijectivité.

Supposons donc d'abord que $u$ est bijectif. Si $H_s$ est \uline{réduit}, on peut évidemment conclure que $u_s$ est un isomorphisme, car $G_s$ s'identifie à un sous-schéma fermé de $H_s$ ayant même ensemble sous-jacent. En particulier, si $k=k(s)$ est de caractéristique nulle, tout groupe algébrique sur $k$ est réduit d'après \uline{Cartier}, et on \deleted{\unsure{a la fin}} a ainsi obtenu:

\uline{Proposition} \deleted{\unsure{6|5}}\inserted{7}.2. Soit $u:G\quad H$ un homomorphisme de préschémas en groupes de présentation finie sur $S$. Supposons que $u$ soit un monomorphisme, $G$ plat sur $S$, et les \deleted{caractéristiques}\inserted{corps} résiduel\deleted{\unsure{les}} de $S$ \deleted{\unsure{tou}} de caractéristique nulle. Alors $u$ est uneimmersion.

Lorsque $k(s)$ est de caractéristique $p\quad0$, nous allons nous borner au cas où $G$ \uline{est commutatif}. Alors (avec les réductions faites) $H$ est également commutatif, car c'est l'adhérence schématique de $H_s$, qui
\mspage{306}
